% Rule of 4 of the brainstem (Gates 2005): cranial nerves by level, midline
% motor nuclei that divide into 12, and the 4 M's (midline) and 4 S's (lateral)
% Art credit: Brain: Servier Medical Art (smart.servier.com), CC BY 4.0 (art/servier-brain.pdf)
\documentclass[11pt, border=2mm]{standalone}
\usepackage{cicm-figure}
\usepackage{array}
\newcommand\organ[2]{\includegraphics[height=#1]{servier-#2.pdf}}
\tikzset{
  midnum/.style = {font=\large\bfseries, text=cicmaccent, fill=cicmaccenttint, inner sep=1.2pt},
  latnum/.style = {font=\bfseries, text=cicmblue},
  level/.style  = {anchor=west, align=left, font=\small},
}
% \level{top y}{bottom y}
\newcommand\level[2]{%
  \fill[cicmbluetint, rounded corners=3mm] (-1.5,#1) rectangle (1.5,#2);
  \fill[cicmaccenttint] (-0.5,#1) rectangle (0.5,#2);
  \draw[rounded corners=3mm] (-1.5,#1) rectangle (1.5,#2);}
\begin{document}
\begin{tikzpicture}
  % Brain for orientation, with the brainstem picked out
  \node (brain) at (-4.6,1.2) {\organ{32mm}{brain}};
  \foreach \a/\b in {{(-4.15,0.95)}/{(-1.5,4.6)}, {(-4.42,0.0)}/{(-1.5,-2.3)}} {
    \draw[white, line width=2.2pt] \a -- \b;
    \draw[thin, dashed, cicmaccent] \a -- \b;}
  \fill[cicmaccent] (-4.15,0.95) circle (0.6mm) (-4.42,0.0) circle (0.6mm);

  % Three levels, midline (M's) tinted crimson, lateral (S's) blue
  \level{4.6}{2.4}
  \level{2.25}{0.05}
  \level{-0.1}{-2.3}
  \draw[thick, cicmaccent] (0,4.75) -- (0,-2.45);

  % Midline motor nuclei: divide evenly into 12
  \node[midnum] at (0,4.0) {3};
  \node[midnum] at (0,3.0) {4};
  \node[midnum] at (0,1.15) {6};
  \node[midnum] at (0,-1.2) {12};

  % Lateral nerves, both sides
  \foreach \s in {-1,1} {
    \node[font=\small, text=cicmgrey] at (\s*1.0,3.5) {--};
    \node[latnum] at (\s*1.0,1.85) {5};
    \node[latnum] at (\s*1.0,1.15) {7};
    \node[latnum] at (\s*1.0,0.45) {8};
    \node[latnum] at (\s*1.0,-0.5) {9};
    \node[latnum] at (\s*1.0,-1.2) {10};
    \node[latnum] at (\s*1.0,-1.9) {11};
  }

  % Level labels
  \node[level] at (1.75,3.5) {\textbf{Midbrain}\\{\footnotesize\color{cicmgrey}CN 1--4 (1, 2 above)}};
  \node[level] at (1.75,1.15) {\textbf{Pons}\\{\footnotesize\color{cicmgrey}CN 5--8}};
  \node[level] at (1.75,-1.2) {\textbf{Medulla}\\{\footnotesize\color{cicmgrey}CN 9--12}};

  % Zone headers
  \node[font=\footnotesize, text=cicmblue] at (-1.0,5.05) {lateral};
  \node[font=\footnotesize, text=cicmaccent] at (0,5.05) {midline};
  \node[font=\footnotesize, text=cicmblue] at (1.0,5.05) {lateral};

  % The 4 M's and the 4 S's
  \fill[cicmaccenttint, rounded corners=0.8mm] (-7.25,-2.9) rectangle (-0.1,-6.0);
  \fill[cicmbluetint, rounded corners=0.8mm] (0.1,-2.9) rectangle (7.25,-6.0);
  \node[anchor=north west, font=\small] at (-6.95,-3.05) {%
    \renewcommand\arraystretch{1.25}\setlength\tabcolsep{4pt}%
    \begin{tabular}{@{}>{\raggedright}p{2.85cm}l@{}}
      \multicolumn{2}{@{}l}{\textbf{\color{cicmaccent}Midline: the 4 M's}}\\
      \textbf{M}otor pathway & contra weakness\\
      \textbf{M}edial lemniscus & contra proprioception\\
      \textbf{M}LF & ipsi INO\\
      \textbf{M}otor nucleus & ipsi CN 3, 4, 6 or 12
    \end{tabular}};
  \node[anchor=north west, font=\small] at (0.4,-3.05) {%
    \renewcommand\arraystretch{1.25}\setlength\tabcolsep{4pt}%
    \begin{tabular}{@{}>{\raggedright}p{2.85cm}l@{}}
      \multicolumn{2}{@{}l}{\textbf{\color{cicmblue}Lateral: the 4 S's}}\\
      \textbf{S}pinocerebellar & ipsi ataxia\\
      \textbf{S}pinothalamic & contra body pain, temp\\
      \textbf{S}ensory nucleus of 5 & ipsi face pain, temp\\
      \textbf{S}ympathetic & ipsi Horner
    \end{tabular}};

  % Take-home
  \node[label, align=center, anchor=north] at (0,-6.2)
    {The cranial nerve gives the level; M's or S's give medial or lateral.\\
     Crossed signs: ipsilateral cranial nerve, contralateral long tracts.};
\end{tikzpicture}
\end{document}
